% Un petit graphe signé, dessiné avec la convention du manuscrit
% (PartIII/ChapII, fig. « Interprétation d'un poids w réel ») et de la
% présentation NIM du 25 juin 2026 :
%   trait plein    = interaction attractive  (W_e > 0)
%   pointillés     = interaction répulsive   (W_e < 0)
%   trait gras     = arête satisfaite par la configuration sigma
% Configuration : a=+, b=+, c=-, d=+, e=-, f=-.
    \begin{tikzpicture}[x=1cm,y=1cm,scale=0.92,every node/.style={scale=0.92}]
      \node[spin] (a) at (0,1.15) {$+$};
      \node[spin] (b) at (1.30,1.55) {$+$};
      \node[spin] (c) at (2.55,1.05) {$-$};
      \node[spin] (d) at (0.30,0.05) {$+$};
      \node[spin] (e) at (1.60,0.35) {$-$};
      \node[spin] (f) at (2.70,-0.15) {$-$};

      % --- arêtes satisfaites (en gras) ---
      \draw[posedge,freeze] (a)--(b);   % attractive, (+,+)
      \draw[posedge,freeze] (a)--(d);   % attractive, (+,+)
      \draw[posedge,freeze] (e)--(f);   % attractive, (-,-)
      \draw[negedge,freeze] (b)--(c);   % répulsive,  (+,-)
      \draw[negedge,freeze] (d)--(e);   % répulsive,  (+,-)

      % --- arêtes non satisfaites (en trait fin) ---
      \draw[posedge,line width=0.6pt] (b)--(e);   % attractive, (+,-)
      \draw[negedge,line width=0.6pt] (c)--(f);   % répulsive,  (-,-)
    \end{tikzpicture}
